\section{Introduction}\label{sec:r37intro}
Dynamic economic decisions depend on continuation values that are costly to
recompute and change when future policies, technologies, or preferences
change. A neural approximation can make repeated evaluation inexpensive, but
neither a small fitting loss nor a successful current decision establishes
that the resulting dynamic policy is accurate. The relevant questions are
which continuation errors affect decisions, how those errors propagate
through the policy's own future, and what work is needed to construct and
verify a neural continuation at the required economic accuracy.

This paper develops Neural Bellman Operators to answer these questions within
a single policy-evaluation and feasible-improvement framework. The original
object is a controlled economy with a policy-specific continuation, including
controlled diffusions and recursive preferences. We retain the distinction
between evaluating a fixed policy and optimizing against that evaluation.
Current actions are varied over their actual feasible set; the future policy
is held fixed while its continuation is evaluated. Boundary, utility-domain,
and comparison conditions belong to the economic problem and are not
replaced by a neural parameterization. The application to which we give a
complete constructive policy certificate is a directly specified recursive
capital economy with Gaussian innovations and vector controls. It is an
application of this framework, not a change in the paper's economic object.

The first contribution identifies the decision-relevant error. An additive
level shift in a continuation does not change an optimizer. The appropriate
local error is therefore an oscillation across feasible actions, together
with any optimization error and any discrepancy between the evaluated and
implemented futures. This distinction matters for both theory and numerical
diagnostics: a method may have a large level error but a small decision loss,
or a small average fitting loss and an inaccurate action. The centered bounds
also explain why a certificate for a current action cannot simply be reused
as a certificate for a changed future.

The second contribution gives a constructive route from trained hidden
weights to full-policy accuracy. In the recursive capital economy, an
independent coefficient comparison applies on all real states and against
all real vector controls. The comparator is every adapted policy with finite
recursive cost, not merely policies in the fitted class. Positive Gaussian
moment margins control the propagation of quadratic and state-independent
errors separately. A primitive envelope allocates policy error before the
neural targets are supplied. Cold construction, a gradient--curvature
criterion, and basin-certified warm construction then provide explicit routes
to the required continuation accuracy. These statements concern an explicitly
trainable square-activation continuation; every entry of its hidden factor
may change. They are not fixed-feature regression guarantees applied to an
unrelated jointly trained network.

The third contribution joins that construction to finite-precision reuse.
A changed continuation may invalidate the old action while leaving the old
factor in an admissible training basin. We derive a two-sided error budget
that pays separately for factor fitting, own-policy evaluation, and the actor
equation. It also checks the magnitude of the fitted coefficient, preventing
an overshooting factor from invalidating the primitive envelope used to prove
the policy bound. Within this budget, a precision-adaptive refresh rule permits
rectangular factors and noncommuting starts. Its relative Gram enclosure
contracts according to $\rho_{j+1}=\rho_j^2$ when the entire implemented
step error receives the stated allowance. A fixed target inverse is constructed
at most once per fit; subsequent applications and checks remain charged.
The resulting theorem specifies both a finite update cap and a sum of stored
fractional precisions. It does not presume that ordinary floating-point
matrix operations have zero error.

These results supply a positive neural construction theorem, not a claim
that neural methods must dominate every structural solution. We distinguish
three statements throughout the paper. A policy certificate establishes the
accuracy of a returned policy. A training theorem establishes that the declared
neural construction can produce an admissible continuation under stated
primitive and initialization conditions. A complete-work comparison establishes
which executed procedure was less costly at a common economic target. None
of these statements alone establishes the other two. In particular, a Gaussian
quadratic economy has a strong structure-exploiting competitor, and we retain
it rather than make the comparison artificially favorable to learning.

The numerical evidence makes these distinctions observable. The earlier
finite-law study establishes valid current-decision certificates and detects
continuation changes, but its adaptive neural rule is more costly than fresh
neural refitting in all eight reported dimension--volume cells; cached
simulation is the least-cost fully certified procedure in those cells. Those
outcomes and the complete original economic applications remain part of the
publication. Subsequent continuous-support, vector-control studies execute
full-policy comparisons, including a 72-service recursive-risk catalogue.
The new frozen comparison isolates construction, precision, and target reuse
on eight exactly specified services. All eight satisfy the same full-policy
tolerance. Precision adaptation reduces storage precision and recorded
complete service time relative to the two inherited neural constructions
in both specified regimes. The structural procedure is faster still.
These are deterministic construction comparisons, not estimates of a
population probability of neural superiority.

The economic interpretation likewise depends on the target. A current
withdrawal response conditional on a specified future policy differs from
a counterfactual in which all future controls are reoptimized. We retain
the former finite-law comparisons with their original interpretation and
use the returned policy's own continuation in every date of the latter
recursive construction. The capital primitives are specified theoretical
economies rather than empirical estimates. The separate applications to
endogenous preferences, recursive utility, temporal selves, and games retain
their own utility domains and deviation or equilibrium conditions. The
capital certificate does not silently discharge those different obligations.

The article proceeds from the economic model and Neural Bellman Operator
to decision errors, full-policy comparison, constructive training, and
precision-adaptive reuse. The technical supplement supplies the proofs and
supporting training criteria. The complete previous article, its supplement,
and all economic applications are preserved as identified companions, with
cross-document references. This organization permits the current argument
to be read in economic and mathematical order without discarding earlier
results or requiring the reader to reconstruct the development from branch
history.
